\section{Empirical Characterization}
\label{sec:empirical}

We evaluated the pipeline across five fully-tested hunt scenarios, spanning four distinct content domains (themes and settings), three delivery formats, and a scale range of 6 to 19 total puzzles. This section reports aggregate pass rates, per-stage timing, failure analysis, and evidence for cross-theme transferability. All measurements are derived from commit timestamps and scenario artifact files in the repository.

%---------------------------------------------------------------------------
\subsection{Scenario Overview}
\label{sec:scenarios}

Table~\ref{tab:scenarios} summarizes the five scenarios that completed the full 11-stage pipeline with verified editorial review. A sixth scenario (Boardgames ``Game Night'') reached Stage 6 before being suspended as a partial workflow test; it is excluded from aggregate pass-rate calculations but its single Stage~6 failure is included in the failure analysis (\S\ref{sec:failures}).

\begin{table}[htbp]
\centering
\caption{Summary of the five fully-tested pipeline scenarios. Pass rate is reported pre- and post-redesign where redesigns occurred. ``Candidates'' is the size of the Stage~3 puzzle pool before selection; ``Final count'' is the number of puzzles shipped.}
\label{tab:scenarios}
\begin{tabular}{@{}llllllll@{}}
\toprule
\textbf{Scenario} & \textbf{Domain} & \textbf{Scale} & \textbf{Delivery} & \textbf{Meta answer} & \textbf{Pool} & \textbf{Final} & \textbf{Pass rate} \\
\midrule
Age of Empires    & Real game (RTS)      & 5+1  & Print/digital    & WOLOLO       & 15 & 6  & 83\%$\to$100\% \\
Grand Larceny     & Noir fiction         & 4+1  & Physical dossier & SLIT         & —  & 6  & 100\%           \\
Ironfall          & Sci-fi RPG (fiction) & 10+3 & Website          & Konami Code  & 18 & 13 & 100\%$^\dagger$ \\
Wavelength        & Music / thematic     & 6+1  & Website + prop   & FADING       & —  & 7  & 100\%           \\
Dead Reckoning    & Sci-fi investigation & 15+4 & Website          & CALIBRATE    & 45 & 19 & 100\% (verified) \\
\midrule
\textbf{Total}    &                      &      &                  &              &    & \textbf{51} & \\
\bottomrule
\multicolumn{8}{@{}l@{}}{\footnotesize $^\dagger$ 100\% post-redesign; 2 puzzles (P01, P02) required extraction redesign at Stage~8.} \\
\multicolumn{8}{@{}l@{}}{\footnotesize ``—'' indicates pool candidates selected within the same session without a discrete count artifact.} \\
\end{tabular}
\end{table}

The five scenarios collectively span four content domains (real-game knowledge, original noir fiction, original sci-fi fiction, and music), three delivery formats (print/digital, physical dossier, and interactive website), and a fourfold scale range (6 to 19 final puzzles). All five ran against the same skill library and reviewer profile set without per-scenario modifications.

%---------------------------------------------------------------------------
\subsection{Pass Rate Analysis}
\label{sec:passrate}

Across the five fully-tested scenarios we track 52 total puzzles (combining the Dead Reckoning verified count with the four live-tested scenarios). Of these, 47 passed first-pass editorial review, yielding an overall first-pass pass rate of \textbf{90.4\%} (47/52). All 5 failures were subsequently recovered through redesign or delivery remediation, bringing the final ship rate to 100\%. These figures measure pass rates against the pipeline's author-constructed AI panel rubric; they do not represent external validation against solver enjoyment or community quality standards.

Figure~\ref{fig:passrate} breaks down the pass/redesign/cut count per scenario. The four scenarios with live solver simulation (Age of Empires, Grand Larceny, Ironfall, Wavelength) total 32 puzzles, of which 29 passed first-pass (90.6\%); 3 required redesign. Dead Reckoning contributes 19 additional verified puzzles at 100\%.

Three observations follow from this distribution. First, pass rate does not degrade with scale: the largest scenario (Dead Reckoning, 19 puzzles) achieved 100\% first-pass, while the smallest (Age of Empires, 6 puzzles) had the only non-100\% first-pass rate among the five. This is consistent with Stage~3 pool ranking becoming more discriminating as the candidate pool grows: the Dead Reckoning pool contained 45 candidates from which 19 were selected, giving the ranking stage more signal to discard weak candidates before authoring began.

Second, pass rate does not correlate with domain type. Fictional worlds (Grand Larceny, Ironfall, Wavelength, Dead Reckoning) performed at least as well as the real-game-knowledge domain (Age of Empires), despite fictional scenarios requiring the pipeline to generate domain facts rather than retrieve them. The world-lock mechanism (\S\ref{sec:system}) appears sufficient to anchor the generation space for fictional domains.

Third, the pattern of which puzzles failed is informative. All three redesigns in the live-tested scenarios occurred at Stage~8 (blind testing and integration), not at earlier stages. Zero failures were introduced after Stage~8 (no puzzle that passed Stage~8 subsequently failed Stage~9 or later on puzzle-quality grounds). This suggests the gates are functioning: failures that the pool ranking and editorial review stages did not catch are surfaced by the blind test, and no quality failures survive to the delivery stage.

%---------------------------------------------------------------------------
\subsection{Stage Timing}
\label{sec:timing}

Precise commit timestamps are available for Scenario~1 (Age of Empires), which was executed in a single uninterrupted session on February~27, 2026. Scenarios 3--5 ran in overlapping parallel sessions on the same day, making per-stage wall-clock attribution ambiguous. Figure~\ref{fig:timing} presents the Scenario~1 timing as a reference baseline.

\begin{table}[htbp]
\centering
\caption{Stage-by-stage elapsed time for Scenario~1 (Age of Empires), derived from commit timestamps. Stage~5 was embedded within Stage~6 and is not separately measured. Stage~4 was resolved in the same session as Stage~3 with no intervening commit.}
\label{tab:timing}
\begin{tabular}{@{}rllr@{}}
\toprule
\textbf{Stage} & \textbf{Name} & \textbf{Commit time} & \textbf{Duration} \\
\midrule
1  & Scope        & 09:16     & —         \\
2  & Structure    & 09:57     & 41 min    \\
3  & Puzzle Pool  & 10:05--10:17 & 8--12 min \\
4  & Assignment   & 10:17     & $\approx$0 min \\
5  & Meta Design  & (embedded in Stage~6) & — \\
6  & Authoring    & 11:06     & 49 min    \\
7  & Editorial    & 11:37     & 31 min    \\
8  & Integration  & 12:01     & 24 min    \\
9  & Live test    & 12:28     & 27 min    \\
10 & Complete     & 12:33     & 5 min     \\
\midrule
\multicolumn{3}{@{}l}{\textbf{Total (Stages 1--10)}} & \textbf{3h 17m} \\
\bottomrule
\end{tabular}
\end{table}

The total elapsed time of 3 hours 17 minutes for a 6-puzzle hunt (5 feeders + meta) from blank scope to ship-ready package establishes a practical throughput baseline. Several features of this timing warrant comment.

Stage~6 (Authoring) at 49 minutes is the longest single stage despite producing the most content (five complete puzzle files). This is because Stage~6 dispatches parallel AI agents, one per puzzle, so effective per-puzzle authoring time is 49 minutes / 5 = approximately 10 minutes per puzzle. The parallelism is the primary throughput mechanism for scale: a 15-puzzle hunt at the same per-puzzle rate would add only marginally to Stage~6 wall time.

Stage~7 (Editorial) at 31 minutes is the most human-intensive review stage: a human reads all AI-generated editorial reports and decides on pass/revise/cut for each puzzle before Stage~8 begins. This human dwell time does not compress with parallelism in the same way Stage~6 does.

Stage~2 (Structure) at 41 minutes reflects the combined cost of world design and round structure generation for this scenario. For fictional scenarios (Ironfall, Grand Larceny), Stage~2 is longer because the world-lock document must be generated from scratch rather than adapted from an existing IP.

These timings are single-run estimates from commit timestamps. Variance across scenarios was not formally measured; the overlapping parallel execution of Scenarios 3--5 precludes clean per-stage attribution for those runs.

%---------------------------------------------------------------------------
\subsection{Failure Analysis}
\label{sec:failures}

Table~\ref{tab:failures} catalogs the five failures observed across all scenarios, including the partial Boardgames scenario. The table codes each failure by stage of detection, failure type, and recovery action.

\begin{table}[htbp]
\centering
\caption{Complete failure inventory across all pipeline runs. ``Stage caught'' identifies where the failure was first detected, not where it originated. GL = Grand Larceny.}
\label{tab:failures}
\begin{tabular}{@{}llllp{4.2cm}@{}}
\toprule
\textbf{\#} & \textbf{Scenario} & \textbf{Puzzle} & \textbf{Stage} & \textbf{Failure type} \& Recovery \\
\midrule
1 & Age of Empires & Puzzle V (Post-Imperial)  & 8 (Blind test)    & Pure computation, no deduction layer; extraction mechanism redesigned \\
2 & Ironfall       & P01 (Bestiary)            & 8 (Integration)   & Answer mismatch (USHER$\to$UMBRA); meta coordination failure; extraction redesigned \\
3 & Ironfall       & P02 (Forge Guide)         & 8 (Integration)   & Answer ambiguity (DHRYY / DRESS / QUELL); extraction redesigned \\
4 & Boardgames     & M2 (Settlers)             & 6 (Authoring)     & Resource math broken, no valid extraction path; not shipped (workflow test) \\
5 & Grand Larceny  & (delivery)                & 9 (Delivery)      & Missing HTML build artifacts; delivery deferred to Stage~11 \\
\bottomrule
\end{tabular}
\end{table}

The five failures divide into four mechanistically distinct categories.

\textbf{Mechanism failure} (failure \#1): Puzzle~V of Age of Empires required solvers to perform integer arithmetic to derive a resource count, with no intermediate deduction step. The Stage~8 blind test identified this as ``pure computation''---the answer was uniquely determined but required calculation rather than insight, violating the puzzle craftsmanship rubric (criterion C3: extraction must reward deduction, not arithmetic). This failure type is in principle detectable at Stage~7 but was missed; Stage~8 provided the corrective catch.

\textbf{Answer coordination failures} (failures \#2 and \#3): Both Ironfall failures arose from a mismatch between feeder-puzzle answers committed during Stage~6 and the meta puzzle's answer network committed during Stage~5. P01 produced the answer USHER, but the meta required UMBRA at that slot; P02 produced an ambiguous extraction that could validly yield three different answers. These failures are the most structurally dangerous in the pipeline: they cannot be detected at Stage~3 (pool ranking occurs before meta design) and are detectable at Stage~7 only if the editorial reviewer has access to the finalized META.md. In this run, the editorial review did not cross-reference META.md, leaving the coordination failure to be caught at Stage~8 integration. This is a process gap, and a design recommendation follows directly: Stage~7 editorial review should explicitly load and verify answer consistency against META.md.

\textbf{Construction error} (failure \#4): The Boardgames meta puzzle M2 (Settlers of Catan variant) was generated with an internal resource-count inconsistency that left no valid extraction path. This failure was caught at Stage~6 during authoring verification---the earliest possible detection point---before any editorial or integration investment was made. The scenario was already a partial workflow test and the puzzle was not shipped.

\textbf{Delivery gap} (failure \#5): Grand Larceny's Stage~9 delivery preparation produced a valid puzzle set but did not generate the HTML build artifacts required for the digital dossier format. This is not a puzzle-quality failure but a pipeline process gap: the \texttt{/hunt site} skill had not yet been invoked, and the delivery checklist did not catch the missing artifacts. Recovery was straightforward (deferred to Stage~11) but the gap illustrates that delivery-format completeness requires explicit checklist enforcement.

The failure distribution by stage (0 at Stage~3, 1 at Stage~6, 3 at Stage~8, 1 at Stage~9, 0 post-Stage~9) supports the primary claim of this paper: the review gates are functioning. Stage~3 pool ranking discards candidates before authoring investment is made; Stage~7 editorial catches mechanism failures; Stage~8 integration catches coordination failures that Stage~7 missed. No puzzle that passed Stage~8 failed on quality grounds at any subsequent stage.

%---------------------------------------------------------------------------
\subsection{Cross-Theme Transferability}
\label{sec:transferability}

A central design goal of the pipeline is transferability across themes and settings: the same skill library and reviewer profile set should run a music-themed hunt and a sci-fi RPG hunt without per-scenario code changes. The evidence supports this claim. All five scenarios share the puzzle hunt format; structural transferability to qualitatively different interactive systems (tabletop campaigns, interactive fiction) remains an open question.

Across the five fully-tested scenarios, the only pipeline artifacts that varied per scenario were the \texttt{SCOPE.md} (hunt identity brief), the \texttt{world/} data files (domain facts and world canon), and the \texttt{ROUNDS.md} (round structure). The skill prompts (\texttt{toolkit/skills/}), reviewer profiles (\texttt{toolkit/profiles/}), rubric criteria, and gate conditions were invariant across all runs.

This invariance is observable directly from the commit history: 142 commits over 2 calendar days include no per-scenario modifications to any skill file. Each scenario's \texttt{CLAUDE.md} imports from \texttt{../../toolkit/} without modification. The toolkit commit messages are domain-agnostic (e.g., \texttt{Add delivery skills: /hunt print, /hunt props, /hunt site}; \texttt{Add /hunt world skill}) and apply equally to all scenarios that were initialized after the commit.

The reviewer profile library did grow during the two-day run: the initial 12-profile panel was expanded to 25 profiles within the first two hours, then to 28, and ultimately to 29 profiles by end of day~2. This growth was driven by perceived coverage gaps (e.g., adding a cryptic-crossword specialist when Ironfall required wordplay-heavy extraction), not by domain-specific customization. The profiles that were added apply across all domains.

%---------------------------------------------------------------------------
\subsection{Pipeline Maturity Growth}
\label{sec:maturity}

The pipeline is not static. The git history records a clear maturation arc across the two-day scenario run: deficiencies identified in early scenarios drove improvements that benefited later scenarios.

The skill library grew from 0 named skills at initialization to 13 skills by the end of the second day. Key inflection points include the addition of \texttt{/hunt world} (commit \texttt{2534b38}, 15:20 on day~1), which was motivated by Ironfall's requirement for a fictional universe design tool not present in the initial toolkit; \texttt{/hunt print}, \texttt{/hunt props}, and \texttt{/hunt site} (commit \texttt{735b379}, 15:58 on day~1), added when scenarios requiring physical and digital delivery revealed gaps in the delivery stage; and \texttt{/hunt resume} and \texttt{/hunt publish} (commit \texttt{3f1936e}, 22:46 on day~1), added to support scenario resumption across sessions and clean distribution packaging.

Bug fixes also propagated forward: two bugs identified in Age of Empires (skills not auto-updating the scenario \texttt{CLAUDE.md} on initialization) were fixed before the Boardgames and Ironfall runs, contributing to the cleaner first-pass behavior observed in the later scenarios.

This maturation pattern is a normal property of pipeline-style system development and is worth naming explicitly: the pass rates reported in \S\ref{sec:passrate} reflect the pipeline in a state of active improvement, not a frozen system. The Age of Empires scenario (first run, 83\% first-pass) and the Dead Reckoning scenario (last run, 100\% verified) bookend the improvement arc. Whether this improvement is attributable to pipeline refinements, to the larger and better-ranked candidate pools in later scenarios, or to accumulated experience within the session is a confound the current study cannot fully disentangle; separating these effects is a direction for future controlled evaluation. Practitioners deploying the pipeline for the first time should expect first-run performance closer to the 83\% observed in Scenario~1 than the 90.4\% aggregate, which blends early and mature runs.
